\section{Introduction}
\begin{figure}[t]
    \centering
    \includegraphics[width=\columnwidth]{figures/Stack.pdf}
    \caption{Citadel's Computing Stack.}
    \label{fig:computing_stack}
    \end{figure}

Citadel proposes a simple and efficient solution to an important problem: ``\emph{How to defend against transient execution side channels for a program that 1) touches secrets, 2) is not written in constant-time, and 3) shares some memory with an attacker}''. 
We call this problem \prob{}. 
It represents a common case for \emph{security-sensitive} programs in the real world which do \emph{not} aim for leakage-free execution times.
For example, consider a private inference service~\cite{ApplePrivateCloud} that needs to isolate the machine learning model and private user requests from the rest of the system, while still accessing the network, a GPU, and communicating with other processes in charge of access control and business logic.

Transient execution side channel (TES) attacks such as Spectre \cite{kocher2018spectre} combine speculative execution and microarchitectural side channels to exfiltrate secrets.
Among the numerous Spectre variants, those exploiting cache side channels are some of the most powerful.
These attacks only require an attacker and the victim to share \emph{any} address space (e.g., through shared memory) to make it possible to construct ``universal read gadgets,'' exposing all the victim's private memory.
They are also particularly challenging to mitigate and circumvent almost all hardware and software isolation mechanisms.
For example, in a scenario like serverless computing, these attacks can compromise the isolation between containers and virtual machines, or enable privilege escalation~\cite{weissman2023microarchitectural}.

Trusted execution environments (TEEs) or secure enclaves \cite{costan2016sanctum, ferraiuolo2017komodo, lee2020keystone,  bahmani2021cure, feng2021scalable, li2022design, cheng2024intel, costan2016intel, mckeen2013sgx, brasser2019sanctuary, sev2020strengthening, trustzone04, pinto2019demystifying} are some of the most ambitious isolation primitives.
These hardware mechanisms aim to completely isolate sensitive processes even in the presence of a privileged software attacker, such as a malicious operating system (OS) or hypervisor (for example, see Figure~\ref{fig:computing_stack}).
However, most of these platforms have been proven vulnerable to TES attacks~\cite{koruyeh2018spectre, chen2019sgxpectre, ryan2019hardware, van2020lvi, van2018foreshadow, van2020sgaxe, gotzfried2017cache, brasser2017software, gras2018translation, evtyushkin2018branchscope, aciiccmez2006predicting, wang2014timing, koruyeh2018spectre}.
On commercial enclave platforms such as Intel SGX~\cite{costan2016intel} or AMD SEV~\cite{sev2020strengthening}, defending against TES is still an open problem, and all responsibility is pushed to program developers.
Countermeasures rely on the placement of fences and special instructions in the binary (see Section \ref{sec:softwaremitigations}), a challenging process that needs to balance preservation of performance and complexity of code analysis, often achieving unclear security guarantees.

One academic work, MI6~\cite{bourgeat2019mi6}, looks at securing enclave programs in the presence of TES through complete microarchitectural isolation.
However, they took the approach to the extreme by disabling shared memory altogether.
This significantly reduces the usability and applicability of enclaves, in contrast to commercial enclave platforms and other isolation primitives such as containers or trusted VMs, which all rely heavily on shared memory to implement many important functionalities, such as input passing, I/O, device virtualization, or micro-services chaining (see Section~\ref{sec:motivsharedmem}).

Beyond enclave-specific solutions, other works have explored software and hardware defense mechanisms against TES attacks (see Section \ref{sec:relatedwork}).
Each of these mechanisms addresses a slightly different trade-off between security, performance, and use cases (see Sections~\ref{sec:beyondsandboxing}).
Within that design space, Citadel and its simple mechanisms stand out as well equipped to efficiently address \prob{}.
\\

\textbf{This paper.}
We start by defining a new speculative security property---\emph{relaxed microarchitectural isolation} (RMI)---that allows \emph{non-constant-time} programs that \emph{touch secrets} to \emph{share memory} (i.e. to address \prob) while providing precise guarantees on information leaked to an attacker.
For a given defense scheme, RMI is achieved if information leakage through the microarchitecture is equivalent to the trace of \emph{shared memory accesses} when executing the program non-speculatively (in-order, one instruction at a time).
We use the formalism from \emph{hardware-software contracts for security}~\cite{guarnieri2020spectector,guarnieri2021hardware, mosier2022axiomatic, cauligi2022sok} to compare and relate RMI with other existing speculative security properties (see Figure~\ref{figure:guarantees}.
We show that RMI is equivalent to a \emph{weak/simple} contract, which makes it possible to secure a very large class of programs despite enforcing a very strong security property under a strong threat model.

We describe RMI in the context of shared memory between an enclave and the OS.
However, our work also applies to \emph{private} shared memory between programs and can help compose code following other non-interference properties such as speculative constant time~\cite{guarnieri2020spectector}.
Our work also extends to other types of isolation mechanisms and TEEs, such as containers or trusted VMs.


To enforce RMI, we design Citadel using two strategies: 1) we use microarchitectural isolation to weaken attacker capabilities and satisfy a weaker leakage model.
2) we use simple hardware mechanisms to constrain speculation in order to enforce simple execution models.
Because we focus on programs that are \emph{not} implemented using constant time and do \emph{not} aim for leakage-free execution times, our threat model allows for some leakage of coarse grain timing information (see Section \ref{sec:execution_time_SC}).
This offers a compelling trade-off between security and enabling \emph{simpler} and \emph{more efficient} mechanisms that can secure a \emph{large class} of enclave programs against Spectre-like attacks.

Starting from a design that enforces strong microarchitectural isolation, we show how to enable shared memory while limiting the attacker observations to the trace of (speculative) accesses to shared memory.
This requires being careful with the core-private state that can now be shared between private and shared memory.
For example, we highlight how an attacker can observe the private state of L1 by exploiting a new side channel that leverages the cache coherence protocol.

Once attacker observations have been restricted, we present two \emph{simple} approaches to limit speculative leakage: (1) {\bf Safe mode} is our default execution mode and prevents speculative shared memory accesses entirely, making it possible to enforce RMI for \emph{any} program, (2) {\bf Burst mode} can be enabled \emph{locally} to allow some speculation on shared memory accesses and improve overall performance. 

The Burst mode microarchitecture enforces a slightly stronger contract compared to Safe mode (more speculative control flows are considered public).
As a result, Burst needs to be supplemented by a static analysis tool that identifies programs that satisfy the correct \emph{relative non-interference} property.
Nevertheless, the property is simple enough that programmers can easily reason about it and can be checked automatically using some simple tool.
We use a contract formalism to show that the resulting hardware-software defense satisfies RMI.
We show how code analysis can be simplified by restricting it to small code snippets, circumventing the usual limitations of automated tools.
Note that code annotation and analysis are only required to improve \emph{performance} and that all programs are secure as is without manual annotation.
Code can be written to use the two modes simultaneously and can be safely composed or reused without the need for additional analysis. 
\\

\textbf{Our prototype.} Citadel, is an out-of-order RISC-V 2-core processor synthesized for an AWS F1 FPGA.
We apply our design principles and implement all mechanisms at the microarchitectural level.
We restrict the attacker's observations through L1-cache bypass and TLB / translation cache tagging.
We also implement a configurable cache partitioning mechanism, improving performance for both enclaves and unprotected processes.
We implement \emph{Safe} mode and \emph{Burst} modes with minimal hardware overhead (<+1\% LUTs and FFs).
To evaluate our prototype, we implement a comprehensive software infrastructure that includes a bootloader, a security monitor, a kernel module, and an end-to-end attestation mechanism.
Citadel boots (untrusted) Linux and enables user-level applications to interact safely with secure enclaves.
For unprotected software, we measure negligible overhead ($\simeq0.1\%$).
To demonstrate the platform's versatility, we port three security-sensitive programs: a lightweight Python runtime, a cryptographic library, and a private ML inference program.
We show that very little implementation effort is required to port the applications to Citadel: Less than $200$ new LOC and minimal automated code analysis for performance overheads under $5\%$.
Figure~\ref{fig:computing_stack} provides an overview of our implementation work.
We open source our hardware and software infrastructure at \url{https://github.com/mit-enclaves/citadel}.

\smallskip\noindent
In summary, our contributions include the following:
\begin{itemize}[leftmargin=*, align=left, itemsep=0ex, topsep=0ex]
\item 
    We present the first microarchitecturally-isolated TEE with Spectre-safe shared memory;
\item 
    We introduce the notion of \emph{relaxed microarchitectural isolation}, a new security property that allows non-constant-time, sensitive programs to share memory while thwarting TES attacks;
\item 
    We demonstrate that microarchitectural isolation and simple hardware mechanisms for controlled speculation make it possible to effectively defend against most TES attacks;
\item 
    We implement and evaluate two simple mechanisms---Safe mode and Burst mode---that achieve RMI for any program while offering the option to balance performance and the need for simple program analysis;
\item 
    We show how to compose code with different speculative security properties and that program analysis can be constrained to small code snippets, circumventing the usual scaling issues of automated tools;
\item
    We use formalism from hardware-software contracts for security to compare RMI with other security property and show that Citadel enforces it;
\item
    We show that our approach allows for minimal hardware overhead (<1\% for Safe and Burst mode) and minimal performance overhead (<5\% on our benchmarks);
\end{itemize}


\bigskip
\noindent\textbf{Disclaimer.} Citadel is a \emph{simple} and \emph{efficient} solution to an important problem.
We do not claim the invention of the simple hardware mechanisms we compose to achieve security.
On the contrary, we aim to highlight that complex mechanisms are not necessary here.
Similarly, most of the insights we present might be considered folklore by some readers.
We show how all these elements are tightly integrated to create a novel point in the design space at the intersection of TEEs and TES defense mechanisms, presenting a compelling trade-off between performance, usability, and security.